\section{Proofs for the boundary results}\label{app:lim}

This appendix proves \cref{thm:lim:ambient,thm:lim:posslp,prop:lim:flow}.

\subsection{Proof of \texorpdfstring{\cref{thm:lim:ambient}}{the ambient-count theorem}}\label{app:lim:ambient}

\paragraph{One block.}
Let $n=m^2$, which is even. Choose $u\in\Q^n$ with $n/2$ entries $1/m$ and
$n/2$ entries $-1/m$, such that $u_{n-1}=u_n=1/m$, and let
$B_+,B_-\subseteq\{1,\ldots,n-2\}$ be the index sets of the remaining positive
and negative entries, so $\abs{B_+}=n/2-2$ and $\abs{B_-}=n/2$. Then
$\norm u=1$. Put $q=e_{n-1}-e_n$, which is orthogonal to $u$, and
$c=\frac34$, $\varepsilon=m^{-3}$,
\[
 A_1=\alpha\bigl[c(uq^{\mathsf T}+qu^{\mathsf T})+qq^{\mathsf T}\bigr],
\]
and let $P_1\subset\R^n$ be the polytope defined by
\[
 x_i\ge0\ (i\le n-2),\qquad\sum_{i\le n-2}x_i=4m,\qquad
 x_{n-1}+x_n=0,\qquad\abs{x_{n-1}-x_n}\le\varepsilon .
\]
The instance of the theorem consists of $k$ copies:
$A=\diag(A_1,\ldots,A_1)$, $T=\diag(u^{\mathsf T},\ldots,u^{\mathsf T})$, and
$P=P_1^k$. Its data have encoding length polynomial in $N$.

\emph{Part \ref{it:lim:amb-instance}.} In the orthonormal basis
$(u,q/\sqrt2)$ of their common range, $A_1$ and $A_1+\alpha uu^{\mathsf T}$
act as
\[
 \alpha\begin{pmatrix}0&c\sqrt2\\c\sqrt2&2\end{pmatrix}
 \qquad\text{and}\qquad
 \alpha\begin{pmatrix}1&c\sqrt2\\c\sqrt2&2\end{pmatrix},
\]
and both vanish on the orthogonal complement. The second matrix has positive
trace and determinant $\alpha^2(2-2c^2)=\frac78\alpha^2>0$, so
$A+\alpha T^{\mathsf T}T\succeq0$. The first has eigenvalues
$\alpha(1\pm\sqrt{1+2c^2})=\alpha(1\pm\sqrt{17/8})$, exactly one of them
negative, and $\nu=\alpha(\sqrt{17/8}-1)>\alpha/4$ because
$\sqrt{17/8}>\frac54$. The rows of $T$ have disjoint supports and unit norm,
so $TT^{\mathsf T}=I_k$. For $x\in P_1$ write $y=u^{\mathsf T}x$ and
$z=q^{\mathsf T}x$. Then $x_{n-1}=z/2$, $x_n=-z/2$, and
$y=m^{-1}(\sum_{B_+}x_i-\sum_{B_-}x_i)$; since the nonnegative $x_i$,
$i\le n-2$, have sum $4m$, the pair $(y,z)$ ranges exactly over
$[-4,4]\times[-\varepsilon,\varepsilon]$. Moreover
$\frac12x^{\mathsf T}A_1x=\alpha(cyz+z^2/2)$.

\paragraph{The auxiliary function of one block.}
Let $\gamma\in\R^n$, $\xi=u^{\mathsf T}\gamma$, $\zeta=\gamma-u\xi$, and
\[
 S=\frac m2\Bigl(\min_{B_+}\gamma_i-\min_{B_-}\gamma_i\Bigr),\qquad
 t=\frac{\gamma_{n-1}-\gamma_n}2 .
\]
We claim that, for a number $K$ depending on $\gamma$ only,
\begin{equation}\label{eq:lim:amb-block}
 V(a)=K+\varphi(a),\qquad
 \varphi(a)=\min_{\abs y\le4,\ \abs z\le\varepsilon}
 \Bigl[G(y,z)+\frac\alpha2\Bigl(a+\frac\xi\alpha-y\Bigr)^2\Bigr],
\end{equation}
with $G(y,z)=\alpha(cyz+z^2/2)+Sy+tz$,
where $V$ is the function \eqref{eq:lim:aux} for one block. Fix $(y,z)$ and
minimize $\zeta^{\mathsf T}x$ over the points of $P_1$ with these
projections. The last two coordinates are determined and contribute
$(\zeta_{n-1}-\zeta_n)z/2=tz$, because $u_{n-1}=u_n$. The
masses $\sum_{B_+}x_i=m(4+y)/2$ and $\sum_{B_-}x_i=m(4-y)/2$ are determined,
and a linear function on a scaled simplex is minimized at a vertex, so the
minimum is $\frac{m(4+y)}2\min_{B_+}\zeta_i+\frac{m(4-y)}2
\min_{B_-}\zeta_i$. Since $\zeta_i=\gamma_i-\xi/m$ on $B_+$ and
$\zeta_i=\gamma_i+\xi/m$ on $B_-$, this equals $K_0+(S-\xi)y$ with
$K_0$ independent of $(y,z)$. Hence
\[
 V(a)=K_0+\xi a+\min_{y,z}\Bigl[\alpha\Bigl(cyz+\frac{z^2}2\Bigr)
 +(S-\xi)y+tz+\frac\alpha2(a-y)^2\Bigr],
\]
and $\frac\alpha2(a-y)^2+\xi(a-y)=\frac\alpha2(a-y+\xi/\alpha)^2-\xi^2/(2\alpha)$
gives \eqref{eq:lim:amb-block} with $K=K_0-\xi^2/(2\alpha)$.

The objective inside \eqref{eq:lim:amb-block} is a quadratic in $(y,z)$ with
Hessian $\alpha\bigl(\begin{smallmatrix}1&c\\c&1\end{smallmatrix}\bigr)\succ0$.
It has a unique minimizer $(y^*(a),z^*(a))$ on the rectangle, and by
Danskin's theorem $\varphi$ is differentiable with
$\varphi'(a)=\alpha(a+\xi/\alpha-y^*(a))$.

\paragraph{A good event.}
Let $\mathcal E_1$ be the event that $\min_{B_+}\gamma_i\le-1+4/n$,
$\min_{B_-}\gamma_i\le-1+4/n$, and $\abs\xi\le2$. We show
$\Prob(\mathcal E_1)>\frac14$ under both laws. Let
$\pi=\Prob\{\gamma_i\le-1+4/n\}$. For the uniform law $\pi=2/n$. For
$U_{1,M}$ the atoms $-1+2j/(M-1)$ in this range are those with
$j\le2(M-1)/n$, at least $2(M-1)/n$ of them, so
$\pi\ge(2/n)(1-1/M)\ge19/(10n)$ because $M\ge32$. In both cases
$\pi\ge19/(10n)$. Using $1-\pi\le e^{-\pi}$, $\abs{B_+}\ge15n/32$ (as
$n\ge64$), and $\abs{B_-}=n/2$,
\[
 \Prob\Bigl\{\min_{B_+}\gamma_i\le-1+\tfrac4n\Bigr\}\ge1-e^{-57/64}>\tfrac7{12},
 \qquad
 \Prob\Bigl\{\min_{B_-}\gamma_i\le-1+\tfrac4n\Bigr\}\ge1-e^{-19/20}>\tfrac35,
\]
because $e^{57/64}>\frac{12}5$ and $e^{19/20}>\frac52$. These two events
depend on disjoint coordinates and are independent. Next,
$\operatorname{Var}\xi=\sum_iu_i^2\operatorname{Var}\gamma_i=\operatorname{Var}
\gamma_1$, which is $\frac13$ for the uniform law and $(M+1)/(3(M-1))\le\frac38$
for the grid law, so $\Prob\{\abs\xi>2\}\le\frac3{32}$ by Chebyshev's
inequality. Therefore $\Prob(\mathcal E_1)\ge\frac7{12}\cdot\frac35-\frac3{32}
=\frac{41}{160}>\frac14$. No independence between $\xi$ and the two minima is
used.

\paragraph{Flatness on the good event.}
On $\mathcal E_1$, both minima lie in $[-1,-1+4/n]$, so
$\abs S\le\frac m2\cdot\frac4n=\frac2m$; also $\abs t\le1$ and
$\abs\xi\le2$.

\emph{Small derivative.} Let $\abs a\le\frac32$. For every
$\abs z\le\varepsilon$, the unconstrained minimizer in $y$ of the objective in
\eqref{eq:lim:amb-block} is $a+\xi/\alpha-(S+\alpha cz)/\alpha$, whose
absolute value is at most $\frac32+2+\frac2m+\varepsilon<4$ because
$\alpha\ge1$. So $y^*(a)$ is this point, and
$\varphi'(a)=S+\alpha cz^*(a)$. Since $\alpha\le m$,
\begin{equation}\label{eq:lim:amb-flat}
 \abs{V'(a)}=\abs{\varphi'(a)}\le\frac2m+\frac34\,m\cdot m^{-3}<\frac3m
 \qquad(\abs a\le\tfrac32).
\end{equation}

\emph{Small gap.} Let $\abs a\le1$. The point $(y,z)=(a+\xi/\alpha,0)$ is
feasible, since $\abs y\le3$, so $\varphi(a)\le G(a+\xi/\alpha,0)\le3\abs S\le
6/m$. Every feasible $(y,z)$ has
$G(y,z)\ge-4\alpha c\varepsilon-4\abs S-\varepsilon\ge-3m^{-2}-8m^{-1}-m^{-3}$,
and $\min\varphi\ge\min G$. Hence
\begin{equation}\label{eq:lim:amb-gap}
 V(a)-\min V=\varphi(a)-\min\varphi<\frac{15}m\qquad(\abs a\le1).
\end{equation}

\paragraph{Products and counts.}
For the $k$-block instance, $A$, $T$, and $P$ split into blocks, so the
minimization in \eqref{eq:lim:aux} splits and
$V(a)=\sum_{l=1}^kV^{(l)}(a_l)$, where $V^{(l)}$ is the one-block
function of the $l$th block of noise coordinates. The block events
$\mathcal E_1^{(l)}$ are independent, so their intersection
$\mathcal E$ has probability greater than $4^{-k}$. The allowance is
$2E_h=\alpha kh^2/4$.

For (b)(i), let $8/\sqrt{\alpha m}\le h\le16/\sqrt{\alpha m}$, so
$2E_h\ge16k/m$. On $\mathcal E$, every $v\in\mathbb G_h\cap[-1,1]^k$ has
$V(v)-\min V=\sum_l(V^{(l)}(v_l)-\min V^{(l)})<15k/m<2E_h$ by
\eqref{eq:lim:amb-gap}. A grid of spacing $h$ has at least
$2/h-1\ge\sqrt{\alpha m}/8-1\ge\sqrt{\alpha m}/16$ nodes in $[-1,1]$, because
$\alpha m\ge256$. So the expected count is at least
$4^{-k}(\sqrt{\alpha m}/16)^k=(\sqrt{\alpha m}/64)^k$.

For (b)(ii), let $16/(\alpha m)\le h\le32/(\alpha m)$. Then $h\le\frac18$, so
$v_l\pm h\in[-\frac32,\frac32]$ for $v_l\in[-1,1]$. On $\mathcal E$,
\eqref{eq:lim:amb-flat} and the mean value theorem give, for every $l$ and
both signs,
\[
 V(v)-V(v\pm he_l)=V^{(l)}(v_l)-V^{(l)}(v_l\pm h)
 \le\frac{3h}m\le\frac{\alpha h^2}4\le2E_h ,
\]
because $h\ge12/(\alpha m)$. The grid has at least
$2/h-1\ge\alpha m/16-1\ge\alpha m/32$ nodes in $[-1,1]$, so the expected count
is at least $4^{-k}(\alpha m/32)^k=(\alpha m/128)^k$.

For (c), $\alpha=1$ and $m=\sqrt{N/k}$ turn the two bounds into
$64^{-k}(N/k)^{k/4}$ and $128^{-k}(N/k)^{k/2}$. The matrix $A$ has $O(N)$
nonzero entries of $O(\log N)$ bits and $P$ has $O(N)$ constraints, so even
a dense encoding has $I\le C_0N^2\log N$ for an absolute constant $C_0$. The
quantities recorded in $R$ do not change with $m$. If either expected count
were at most $f(k,R)I^c$, then fixing $k>8c$ and letting $m$ grow would
contradict $64^{-k}(N/k)^{k/4}\le f(k,R)(C_0N^2\log N)^c$.

Finally, a search grid of the low-rank route qualifies: it divides an
auxiliary box that contains $[-\frac32,\frac32]$ in every coordinate by
powers of two, and its meshes $w2^{-j}$ meet every interval $[t,2t]$ with
$t$ at most the box width $w$, in particular both ranges of $h$ in (b). The
ambient diameter of $P_1$ is at least $4\sqrt2\,m$, since $P_1$ contains
$4me_i$ for every $i\le n-2$; the theorem therefore does not contradict
bounds that grow with the ambient diameter.\qed

\subsection{Proof of \texorpdfstring{\cref{thm:lim:posslp}}{the point-output theorem}}\label{app:lim:posslp}

The constructions of parts (a) and (b), and the deterministic consequence,
also appear in the unpublished companion
manuscript~\citep{companion-exact-arithmetic}; they are reproduced here so
that the proof is self-contained. Both use the same two tools.

\begin{lemma}[A sign test]\label{lem:lim:signs}
Let $\Phi$ be a polynomial on a box $Z$ with $\nabla^2\Phi\succeq\kappa I$ on
$Z$ for some $\kappa\in(0,1]$, let $z^0$ be its unique minimizer on $Z$, and
let $\zeta$ be an affine function with $\norm{\nabla\zeta}\le1$. Put
$\varepsilon=\kappa/4$ and, on $Z\times[0,1]$,
\[
 H(z,\theta)=\Phi(z)+\theta^2+\varepsilon\,\theta\,\zeta(z).
\]
Then $\nabla^2H\succeq\frac34\kappa I$, $H$ has a unique minimizer
$(z^H,\theta^H)$, and $\theta^H>0$ if and only if $\zeta(z^0)<0$.
\end{lemma}

\begin{proof}
The Hessian of $\Phi(z)+\theta^2$ is at least $\kappa I$, and the bilinear
term has Hessian of operator norm $\varepsilon\norm{\nabla\zeta}\le\kappa/4$.
This gives the modulus and uniqueness on the compact convex box. On the face
$\theta=0$, $H$ equals $\Phi$, whose unique minimizer is $z^0$. At
$(z^0,0)$ the $z$-gradient of $H$ is $\nabla\Phi(z^0)$, which satisfies the
first-order conditions of $\Phi$ on $Z$, and $\partial_\theta H=\varepsilon
\zeta(z^0)$. If $\zeta(z^0)\ge0$, the point $(z^0,0)$ satisfies the
first-order conditions of the convex function $H$ on its box, so it is the
minimizer and $\theta^H=0$. If $\zeta(z^0)<0$, increasing $\theta$ decreases
$H$, so $(z^0,0)$ is not the minimizer; no other point of the face
$\theta=0$ is either, since a minimizer on that face minimizes $\Phi$. Hence
$\theta^H>0$.
\end{proof}

\begin{lemma}[A convex amplifier]\label{lem:lim:amplifier}
Let $H_\pm$ be polynomials on boxes $Z_\pm\times[0,1]$, with last coordinates
$\theta_\pm$, whose Hessians are at least $\mu_0I$, and whose unique
minimizers have exactly one of $\theta_+^H,\theta_-^H$ positive. Put
$\eta=\mu_0/12$ and, with $y\in[0,1]$,
\[
 \Phi=H_+(z_+,\theta_+)+H_-(z_-,\theta_-)
 +\eta\bigl[\theta_+^2(y-1)^2+\theta_-^2y^2\bigr].
\]
Then $\Phi$ is convex on its box and has a unique minimizer, whose
$y$-coordinate is $1$ if $\theta_+^H>0$ and $0$ if $\theta_-^H>0$.
\end{lemma}

\begin{proof}
For a constant $c$ and a direction $(\pi,\varsigma)$ in $(\theta,y)$, direct
differentiation gives
\[
 (\pi,\varsigma)^{\mathsf T}\nabla^2\bigl[\theta^2(y-c)^2\bigr](\pi,\varsigma)
 =2\bigl(\theta\varsigma+2(y-c)\pi\bigr)^2-6(y-c)^2\pi^2 .
\]
With $c=1$ for $\theta_+$, $c=0$ for $\theta_-$, and $\abs{y-c}\le1$, the
amplifier contributes at least $-6\eta(\pi_+^2+\pi_-^2)$, where $\pi_\pm$ are
the components of a direction along $\theta_\pm$. The base Hessian
contributes at least $\mu_0$ times the squared norm of the base part of the
direction. So the Hessian quadratic form is at least $\frac12\mu_0$ times
that squared norm, and $\Phi$ is convex. No division by $\theta_\pm$ occurs,
so this includes points where both amplitudes vanish.

The amplifier is nonnegative, so $\Phi\ge\min H_++\min H_-$. If
$\theta_+^H>0=\theta_-^H$, the two base minimizers with $y=1$ make the
amplifier vanish, so the bound is attained; every minimizer must attain both
base minima, which fixes all base coordinates, and must make the amplifier
vanish, which forces $y=1$. The case $\theta_-^H>0$ is symmetric and forces
$y=0$.
\end{proof}

\paragraph{Part (a): Square Root Sum.}
An instance consists of positive integers $a_1,\ldots,a_n$ and an integer
$B$. If $B<0$, the answer is no. Compute powers of two $R_i$ with
$R_i^2\le a_i<4R_i^2$, and put $c_i=a_i/R_i^2\in[1,4)$, $R=\max_iR_i$,
$w_i=R_i/R\in(0,1]$, and let $K$ be the least power of two with
$K\ge\max\{2,n\}$. Take a complete binary tree with $K$ leaf slots. Each real
leaf $i\le n$ has a variable $u_i\in[1,2]$ and signal $w_iu_i$; padded leaves
have signal $0$. Each internal node $\nu$ has a variable $v_\nu\in[0,2]$,
its own signal $v_\nu$, and residual
$\rho_\nu=v_\nu-\frac12(\text{left signal}+\text{right signal})$. Let $s$ be
the root variable, $Y$ the box of all these variables, and
\[
 F_{\rm tree}(u,v)=\sum_{i=1}^n\Bigl(\frac{u_i^3}3-c_iu_i\Bigr)+\sum_\nu\rho_\nu^2 .
\]

\begin{lemma}[Averaging tree]\label{lem:lim:tree}
$\nabla^2F_{\rm tree}\succeq\frac18I$ on $Y$. The unique minimizer has
$u_i=\sqrt{c_i}$, all residuals zero, and root value
$s^0=\sum_i\sqrt{a_i}/(KR)\in[0,2]$. The bags consisting of an internal
variable and its variable children, joined along the tree, form a tree
decomposition of bag size at most three that covers every term.
\end{lemma}

\begin{proof}
The leaf term has derivative $u_i^2-c_i$ and is minimized on $[1,2]$ at
$\sqrt{c_i}\in[1,2)$. Setting each internal variable to the average of its
children's signals makes every residual zero and keeps it in $[0,2]$; this
point minimizes every summand, and its root is the average of the $K$ leaf
signals, $K^{-1}\sum_iw_i\sqrt{c_i}=\sum_i\sqrt{a_i}/(KR)$. Write the
residuals and the leaf variables of a direction $d$ as $(I-P)d$, where a
leaf row of $P$ is zero and the row of an internal node holds the
coefficients $\frac12w_i$ or $\frac12$ of its variable children. Every
variable is a child of at most one node, so distinct rows of $P$ have
disjoint supports, $PP^{\mathsf T}$ is diagonal with entries at most
$\frac12$, and $\norm P\le1/\sqrt2$. Since $2u_i\ge2$,
\[
 d^{\mathsf T}\nabla^2F_{\rm tree}\,d\ge2\norm{(I-P)d}^2
 \ge2\bigl(1-1/\sqrt2\bigr)^2\norm d^2\ge\tfrac18\norm d^2 .
\]
This also proves uniqueness. A leaf variable occurs only in its parent's
bag, which covers its unary term and its parent's residual; adjacent bags
share one internal variable.
\end{proof}

\begin{lemma}[Integral radical sums]\label{lem:lim:trace}
If positive integers $a_1,\ldots,a_n$ satisfy $\sum_i\sqrt{a_i}\in\Z$, then
every $a_i$ is a perfect square.
\end{lemma}

\begin{proof}
Let $\mathcal F=\Q(\sqrt{a_1},\ldots,\sqrt{a_n})$ and
$\operatorname{Tr}=\operatorname{Tr}_{\mathcal F/\Q}$. If $a_i$ is not a
square, $\Q(\sqrt{a_i})$ has degree two and the trace of $\sqrt{a_i}$ from it
to $\Q$ is zero, so transitivity gives $\operatorname{Tr}(\sqrt{a_i})=0$. A
rational $z$ has $\operatorname{Tr}(z)=[\mathcal F:\Q]z$. Applying the trace
to $\sum_i\sqrt{a_i}=B'\in\Z$ shows that the sum over the square radicands
alone equals $B'$, so the sum of the positive numbers $\sqrt{a_i}$ over the
nonsquares is zero, and there are no nonsquares.
\end{proof}

The field $\mathcal F$ is used only in the proof; the reduction tests
squareness by integer square roots. If every $a_i$ is a square, decide the
instance by integer arithmetic. If $B\ge2KR$, the answer is yes, because
$s^0\le2$. Otherwise $b=B/(KR)\in[0,2)$, and $s^0\ne b$ by
\cref{lem:lim:trace}. Take two copies of $F_{\rm tree}$ with root variables
$s_\pm$, apply \cref{lem:lim:signs} with $\kappa=\frac18$ to the first copy
with $\zeta_+=b-s_+$ and to the second with $\zeta_-=s_--b$, and apply
\cref{lem:lim:amplifier} with $\mu_0=\frac3{32}$. Then $\theta_+^H>0$ exactly
when $s^0>b$ and $\theta_-^H>0$ exactly when $s^0<b$, so the resulting
quartic $\Phi$ is convex on its box and its unique minimizer has
$y$-coordinate $1$ if $\sum_i\sqrt{a_i}>B$ and $0$ if $\sum_i\sqrt{a_i}<B$.

For the structure, take the two decompositions of \cref{lem:lim:tree}, add
the bags $\{s_\pm,\theta_\pm\}$ at their roots and the bags
$\{\theta_\pm,y\}$ attached to them, and join the two new bags; this is a
tree decomposition of bag size three that covers every term, so the
interaction graph has treewidth at most two. The diagonal Hessian entries
are at most $\frac92$ for leaf variables, $\frac52$ for internal variables,
$2+2\eta$ for $\theta_\pm$, and $4\eta$ for $y$, hence at most $5$. All
coordinate intervals have width at most two. The affine maps onto $[0,1]$
preserve degree and scopes and multiply diagonal Hessian entries by at most
four, giving the bound $20$. Every term has degree at most four, at most
three variables, and coefficients bounded by an absolute constant after the
substitution, and every variable lies in at most four terms, so every
nonconstant monomial of the expanded polynomial has a coefficient bounded by
an absolute constant. The constant term need not be bounded, and deleting it
does not change the minimizer. The construction takes polynomial time.

\paragraph{Part (b): $\PosSLP$.}
A straight-line program is a sequence of integers $A_1,\ldots,A_s$ in which
each $A_j$ is $0$, $1$, or $A_l\circ A_{l'}$ with $l,l'<j$ and
$\circ\in\{+,-,\times\}$; its output is $A=A_s$. The integers can have
exponentially many bits. The construction never computes them; it writes
only bounded local formulas.

\begin{lemma}[Bounded pair encoding]\label{lem:lim:pairs}
Define rational pairs $(u_j,v_j)$ by $(0,\frac14)$ for the constant $0$,
$(\frac14,\frac14)$ for the constant $1$, and, for $A_j=A_l\pm A_{l'}$ and
$A_j=A_lA_{l'}$ respectively,
\[
 u_j=\frac{u_lv_{l'}\pm u_{l'}v_l}4,\quad v_j=\frac{v_lv_{l'}}4,\qquad
 u_j=\frac{u_lu_{l'}}4,\quad v_j=\frac{v_lv_{l'}}4 .
\]
Then $v_j>0$, $u_j=A_jv_j$, and $\abs{u_j},\abs{v_j}\le\frac14$ for every $j$.
The number $t=(2u_s-v_s)/4=v_s(2A-1)/4$ is nonzero, it is positive if and
only if $A>0$, and $\abs t\le\frac3{16}$.
\end{lemma}

\begin{proof}
By induction, $v_j$ is a product of positive numbers, and the ratio
$u_j/v_j$ equals $u_l/v_l\pm u_{l'}/v_{l'}$ or $(u_l/v_l)(u_{l'}/v_{l'})$,
which is $A_j$. If all earlier coordinates lie in $[-\frac14,\frac14]$, the
new numerators have absolute value at most $2\cdot\frac1{16}\cdot\frac14=
\frac1{32}$ or $\frac1{64}$. Since $A$ is an integer, $2A-1\ne0$, and $v_s>0$.
\end{proof}

List the coordinates $u_1,v_1,\ldots,u_s,v_s,t$ as $x_1,\ldots,x_N$, $N=2s+1$.
Each is given by a rule $x_i=p_i(x_1,\ldots,x_{i-1})$, where $p_i$ is a
constant, one of the bilinear formulas of \cref{lem:lim:pairs}, or the
affine formula for $t$. On the box $\mathcal B=[-\frac14,\frac14]^N$ these
rules satisfy
\begin{equation}\label{eq:lim:gate-bounds}
 \abs{p_i}\le\tfrac14,\qquad\norm{\nabla p_i}_1\le1,\qquad
 \norm{\nabla^2p_i}_2\le1 .
\end{equation}
Indeed, each bilinear rule has at most four gradient entries of absolute
value at most $\frac1{16}$, or fewer and larger entries when $l=l'$, with
$\norm{\nabla p_i}_1\le\frac14$; its Hessian has at most two nonzero entries
in each row, of absolute value at most $\frac14$, or a single entry at most
$\frac12$ when $l=l'$; and the rule for $t$ is affine with
$\norm{\nabla p_N}_1=\frac34$. The exact assignment $x^*$, defined by
$x_i^*=p_i(x_{<i}^*)$, lies in $\mathcal B$ by \cref{lem:lim:pairs}, and
$x_N^*=t^*$ has the sign of $A-\frac12$.

\begin{lemma}[Weighted gate objective]\label{lem:lim:gates}
Let $G(x)=\sum_{i=1}^N64^{N-i}\bigl(x_i-p_i(x)\bigr)^2$. Then
$\nabla^2G\succeq I$ on $\mathcal B$, and $x^*$ is the unique minimizer of $G$
on $\mathcal B$.
\end{lemma}

\begin{proof}
Write $a_i=64^{N-i}$, $\rho_i=x_i-p_i(x)$, $D=\diag(\sqrt{a_i})$, and let $L$
be the strictly lower-triangular matrix with $L_{ij}=\partial_jp_i$. For a
direction $h$,
\[
 h^{\mathsf T}\nabla^2Gh=2\norm{D(I-L)h}^2-2\sum_ia_i\rho_i\,h^{\mathsf T}\nabla^2p_ih .
\]
Put $\widehat E=DLD^{-1}$, so $\widehat E_{ij}=8^{j-i}L_{ij}$ for $j<i$. By
\eqref{eq:lim:gate-bounds} every row of $\widehat E$ has absolute sum at
most $\frac18$, and every column has absolute sum at most
$\sum_{l\ge1}8^{-l}=\frac17$. Hence
$\norm{\widehat E}_2\le(\frac18\cdot\frac17)^{1/2}<\frac14$ and
\[
 2\norm{D(I-L)h}^2=2\norm{(I-\widehat E)Dh}^2\ge\tfrac98\norm{Dh}^2 .
\]
On $\mathcal B$, $\abs{\rho_i}\le\frac12$, and $p_i$ depends only on
$x_{<i}$, so $\abs{h^{\mathsf T}\nabla^2p_ih}\le\norm{h_{<i}}^2$. Therefore the
second term is at least
\[
 -\sum_ia_i\norm{h_{<i}}^2=-\sum_jh_j^2\sum_{i>j}a_i\ge-\frac1{63}\sum_ja_jh_j^2 .
\]
Altogether $h^{\mathsf T}\nabla^2Gh\ge(\frac98-\frac1{63})\norm{Dh}^2\ge\norm h^2$,
since every $a_j\ge1$. Because the rules are triangular, $G\ge0$ vanishes
exactly at $x^*$, which is therefore the unique minimizer.
\end{proof}

Take two copies of $G$ with outputs $t_\pm$, apply \cref{lem:lim:signs} with
$\kappa=1$, $\zeta_+=-t_+$, and $\zeta_-=t_-$, and apply
\cref{lem:lim:amplifier} with $\mu_0=\frac34$. Since $t^*\ne0$,
$\theta_+^H>0$ exactly when $A>0$ and $\theta_-^H>0$ exactly when $A\le0$.
The resulting quartic $\Phi$ is convex on its rational box, and its unique
minimizer has $y$-coordinate $1$ if $A>0$ and $0$ if $A\le0$. Every residual
$x_i-p_i(x)$ has at most three monomials and every weight has $O(N)$ bits,
so $\Phi$ is an explicit rational quartic with $O(N)$ monomials, computed in
time polynomial in the size of the program.

\paragraph{Part (c).}
For $F_\gamma$, the core and residual terms are separate. The core term
$(v-\frac12)^2+\gamma v$ equals $(v-v_\gamma)^2$ plus a constant, so the
core optimizer is $v_\gamma=(1-\gamma)/2\in[\frac14,\frac34]$, the projected
growth is one, and $\partial_{vv}F_\gamma=2$. The residual objective $\Phi$
does not depend on $\gamma$, so the unique optimizer is $(v_\gamma,\hat x)$
on every draw. Given an instance of either problem, construct $\Phi$,
compute the law, draw $\gamma$, run the algorithm, and compute a point within
distance $1/4$ of $(v_\gamma,\hat x)$. Its $y$-coordinate is within $1/4$ of
$\hat x_y$, so comparing it with $\frac12$ decides the instance. The answer
is correct on every draw, and the expected running time is polynomial. For
the deterministic statement apply the algorithm to $\Phi$ directly. Finally,
Square Root Sum is decidable in polynomial time with a $\PosSLP$
oracle~\citep{allender2009-on-the-complexity-of-numerical}, which gives the
consequences for Square Root Sum also through part (b).\qed

\subsection{Proof of \texorpdfstring{\cref{prop:lim:flow}}{the flow-boundary proposition}}\label{app:lim:flow}

The atoms of $U_{1,M}$ are $-1+2j/(M-1)$, $0\le j<M$. Such an atom is at
least $\frac12$ exactly when $j\ge3(M-1)/4$. For a power of two $M\ge4$,
$3M/4$ is an integer and $3M/4-1<3(M-1)/4\le3M/4$, so the condition is
$j\ge3M/4$, which holds for $M/4$ indices. Each coordinate therefore has
probability $\frac14$, and by independence $\Pr(\mathcal E)=\frac1{16}$. The
second derivatives of $F_\gamma$ in the core are $\partial_{v_iv_i}=1$, so
$L=1$ is valid.

\emph{Part \ref{it:lim:flow-growth}.} Minimizing $v_1z_1+v_2z_2$ over the
two labels gives $\min\{v_1,v_2\}$, so $V_\gamma$ has the stated form. On
$\mathcal E$ each $\gamma_i>0$, and $v_i\ge v_i^2$ on $[0,1]$, so
$\gamma^{\mathsf T}v\ge\min_i\gamma_i\norm v^2$; with $\min\{v_1,v_2\}\ge0$
this gives $V_\gamma(v)\ge(\frac12+\min_i\gamma_i)\norm v^2\ge\norm v^2$.
Since $V_\gamma(0)=0$, the origin is the unique optimal core, and both
labels have cost zero there.

\emph{Part \ref{it:lim:flow-box}.} At $(a,0)$ with $a>0$ the label $(1,0)$
costs $a$ and $(0,1)$ costs $0$; at $(0,b)$ with $b>0$ the opposite holds.
The noise term is common to both labels and cancels from these comparisons.
So each label is strictly worse than the other at some point of the box,
and no valid certificate can assert that one label is optimal throughout.

\emph{Part \ref{it:lim:flow-search}.} The cell $[0,h]^2$ has the corner $0$
with value $V_\gamma(0)=\min V_\gamma$, and the incumbent $U_j$ is a
feasible value, hence at least $\min V_\gamma$. So the lower bound of the
cell, at most $V_\gamma(0)-h^2/4$, does not exceed $U_j$, and the cell is
retained whenever its parent is. At level $0$ the whole core box is
retained, and induction gives the claim at every level.

\emph{The chain.} With $m-1$ further stages of two zero-cost parallel arcs,
each feasible flow chooses one arc per stage, so there are $2^m$ flows. The
conditional value is unchanged, so parts
\ref{it:lim:flow-growth}--\ref{it:lim:flow-search} hold, and on
$\mathcal E$ every flow is optimal at the origin. In every box
$[0,a]\times[0,b]$ the optimal flows at $(a,0)$ and at $(0,b)$ use
different first-stage arcs, so no single flow is certified on the retained
hull at any level. A strategy that enumerates all flows when this
certificate fails does at least $2^m$ work on $\mathcal E$, and its
expected work is at least $2^m/16$. The network has $2m$ unit-capacity arcs
and $m+1$ nodes, so its encoding has length $O(m\log m)$.\qed
